\documentclass[12pt]{article}
\usepackage[utf8]{inputenc}
\usepackage[spanish]{babel}
\usepackage{amsmath, amssymb, amsfonts}
\usepackage{graphicx}
\usepackage{hyperref}
\usepackage{geometry}
\geometry{a4paper, margin=1in}

\title{Proyecto 2: Valoración de Opciones Europeas}

\begin{document}

\maketitle

\section{Introducción y Motivación}
El presente documento tiene como objetivo principal explorar, de forma secuencial y estructurada, los métodos matemáticos y computacionales para la valoración de opciones financieras y la gestión del riesgo (cobertura). 

Partiendo desde la concepción del mercado bajo el modelo del movimiento browniano geométrico, se avanza hacia la derivación analítica de la ecuación de Black-Scholes para un derivado exótico conocido como opción ``Call Ventana''. Finalmente, el proyecto converge en la aplicación práctica de esta teoría mediante el diseño de una estrategia de portafolio, demostrando con simulaciones cómo la elección de este instrumento acotado permite asegurar posiciones y alterar asimétricamente la exposición al riesgo. En conjunto, estos elementos proveen un marco integral que conecta la teoría estocástica pura con la toma de decisiones para la gestión táctica de inversiones.

\section{Simulación del Camino Aleatorio del Precio de un Activo}

\subsection{Descripción del Modelo}
Considere un activo cuyo precio sigue un camino aleatorio dado por la siguiente ecuación diferencial estocástica:
\begin{equation}
dS = \mu S dt + \sigma S dW
\end{equation}
donde:
\begin{itemize}
    \item $S$ es el precio del activo.
    \item $\mu$ es la tasa de rendimiento esperada (tendencia o \textit{drift}).
    \item $\sigma$ es la volatilidad del activo.
    \item $dW$ es el incremento de un proceso de Wiener estándar, es decir, una variable aleatoria con distribución normal de media nula y varianza $dt$.
\end{itemize}

Este modelo corresponde al movimiento browniano geométrico, utilizado en el modelo de Black-Scholes para describir la evolución del precio de un activo subyacente.

\subsection{Metodología de Simulación}

Para simular computacionalmente este proceso, se debe discretizar la ecuación diferencial en el tiempo utilizando el método de Euler-Maruyama. Sea $[t_0, T]$ el intervalo de tiempo a simular. Se divide este intervalo en $N$ pasos de tiempo, cada uno con un incremento constante:
\begin{equation}
\Delta t = \frac{T - t_0}{N}
\end{equation}

La sucesión de tiempos de simulación está dada por $t_{i+1} = t_i + \Delta t$, para $i = 0, 1, \dots, N-1$.

La fórmula recursiva para el precio simulado en el tiempo $t_{i+1}$ a partir del tiempo $t_i$ se expresa como:
\begin{equation}
S_{i+1} = S_i + \mu S_i \Delta t + \sigma S_i \Delta W_i
\end{equation}
donde $\Delta W_i = W_{i+1} - W_i$ es el incremento del proceso de Wiener en el paso de tiempo $i$. Se sabe que este incremento se distribuye normalmente con media 0 y varianza $\Delta t$. Por lo tanto, podemos reescribir el incremento como:
\begin{equation}
\Delta W_i = \sqrt{\Delta t} Z_i
\end{equation}
donde $Z_i$ es una variable aleatoria generada con distribución normal estándar ($Z_i \sim \mathcal{N}(0,1)$).

Sustituyendo esto en la fórmula recursiva, obtenemos la ecuación a implementar en el algoritmo:
\begin{equation}
S_{i+1} = S_i + \mu S_i \Delta t + \sigma S_i \sqrt{\Delta t} Z_i
\end{equation}
para $i = 0, 1, \dots, N-1$.

\subsection{Algoritmo Computacional}

El algoritmo paso a paso para simular una trayectoria del precio del activo es el siguiente:

\begin{enumerate}
    \item \textbf{Inicialización de parámetros:}
    \begin{itemize}
        \item Definir el precio inicial $S_0$.
        \item Definir el tiempo inicial $t_0$ y el tiempo final $T$ de la simulación.
        \item Definir el número de pasos de tiempo $N$.
        \item Ingresar la tasa de interés o rendimiento esperado $\mu$ y la volatilidad $\sigma$.
    \end{itemize}
    \item \textbf{Cálculo del paso de tiempo:} Calcular el incremento temporal finito $\Delta t = (T - t_0) / N$.
    \item \textbf{Inicialización de variables de estado:} Asignar $S(0) = S_0$ y $t(0) = t_0$.
    \item \textbf{Bucle de simulación:} Para $i = 0$ hasta $N-1$, iterar los siguientes pasos:
    \begin{enumerate}
        \item Calcular el siguiente paso temporal: $t_{i+1} = t_i + \Delta t$.
        \item Generar un número aleatorio pseudocongruencial $Z_i$ que siga una distribución normal estándar $\mathcal{N}(0,1)$.
        \item Calcular el siguiente valor simulado del activo:
        \begin{equation*}
        S_{i+1} = S_i + \mu S_i \Delta t + \sigma S_i \sqrt{\Delta t} Z_i
        \end{equation*}
    \end{enumerate}
    \item \textbf{Retorno de resultados:} Retornar la serie de tiempos $\{t_0, t_1, \dots, t_N\}$ y la trayectoria simulada de precios $\{S_0, S_1, \dots, S_N\}$.
\end{enumerate}

Este algoritmo permite obtener múltiples trayectorias (escenarios) independientes si se ejecuta repetidas veces con los mismos parámetros pero generando distintas secuencias de variables aleatorias $Z_i$. Al encapsular esto en una función computacional, el usuario podrá modificar libremente los parámetros $(\mu, \sigma, S_0, T, N)$ y observar el comportamiento de las distribuciones de los caminos estocásticos resultantes.

\subsection{Resultados de la Simulación y Análisis}

Para ilustrar el comportamiento del modelo estocástico propuesto y su sensibilidad ante diferentes escenarios económicos, se ha implementado el algoritmo computacional en lenguaje Python y se han ejecutado múltiples simulaciones. Se evaluó el modelo considerando una matriz de variaciones de $3 \times 3$, analizando el efecto de tres niveles de la tasa de interés ($\mu$: Baja al 2\%, Media al 8\% y Alta al 20\%) frente a tres niveles de volatilidad ($\sigma$: Baja al 5\%, Media al 20\% y Alta al 40\%). Se generaron 10 trayectorias (caminos aleatorios) por cada escenario asumiendo un precio inicial $S_0 = 100$ durante un periodo $T=1$ año ($N=252$ días).

En la Figura \ref{fig:simulacion} se presentan las 9 gráficas correspondientes.

\begin{figure}[htpb]
    \centering
    \includegraphics[width=\textwidth]{simulacion_3x3.png}
    \caption{Simulación del precio del activo bajo distintos escenarios de tasa de interés ($\mu$) y volatilidad ($\sigma$). La línea punteada indica el precio inicial $S_0$.}
    \label{fig:simulacion}
\end{figure}

\subsubsection{Análisis y Conclusión}

A partir de la inspección de la Figura \ref{fig:simulacion}, se pueden extraer las siguientes conclusiones fundamentales sobre la dinámica del modelo:

\begin{itemize}
    \item \textbf{Efecto de la Tendencia ($\mu$):} Al observar las gráficas moviéndose de arriba hacia abajo (incremento en $\mu$), se hace evidente que la tasa de interés o rendimiento esperado dicta la \textit{tendencia general} del camino. Con una $\mu$ alta ($20\%$), el "drift" positivo domina fuertemente, y la mayoría de los escenarios resultan en precios finales $S_T$ significativamente mayores que el precio inicial $S_0$. 
    \item \textbf{Efecto de la Volatilidad ($\sigma$):} Al observar las gráficas moviéndose de izquierda a derecha (incremento en $\sigma$), la amplitud de las fluctuaciones estocásticas se magnifica. En escenarios de volatilidad baja ($5\%$), las trayectorias se agrupan de forma ceñida alrededor de la tendencia determinista. Por el contrario, en escenarios de volatilidad alta ($40\%$), el componente aleatorio del término de difusión de Wiener ($dW$) asume el protagonismo; las fluctuaciones de precios son tan agresivas que los caminos divergen ampliamente entre sí. En estos casos altamente volátiles, incluso cuando hay una tasa de interés $\mu$ positiva y alta, existe un riesgo considerable de que el precio termine por debajo del valor inicial, haciendo que la predicción del valor final sea altamente incierta.
\end{itemize}

\textbf{Conclusión:} El modelo de Black-Scholes a través del movimiento browniano geométrico capta con fidelidad la relación riesgo-retorno del mercado financiero. La variable $\mu$ proyecta hacia dónde debería tender a largo plazo el valor del activo, mientras que $\sigma$ representa la incertidumbre y el grado de dispersión esperado alrededor de esa tendencia. Por lo tanto, para una correcta gestión de portafolios e instrumentos financieros como opciones (ej. \textit{Call Ventana}), no basta con pronosticar rendimientos; es mandatorio calibrar rigurosamente la volatilidad, ya que ésta puede anular fácilmente la tendencia dominante en escenarios de alta incertidumbre.

\section{Valoración de Opción ``Call Ventana''}

\subsection{Ecuación de Black-Scholes para el Modelo Estocástico}
La función $C = C(S,t)$ representa el valor de la opción ``Call Ventana''. Suponiendo un mercado sin oportunidades de arbitraje y donde el activo subyacente sigue el movimiento browniano geométrico planteado, el valor de esta opción satisface la ecuación diferencial parcial de Black-Scholes:
\begin{equation}
\frac{\partial C}{\partial t} + \frac{\sigma^2 S^2}{2} \frac{\partial^2 C}{\partial S^2} + r S \frac{\partial C}{\partial S} - r C = 0
\end{equation}
donde $r$ es la tasa libre de riesgo, $\sigma$ es la volatilidad del activo y $t \in [0, T]$.

\subsection{Condiciones de Frontera}
Una opción ``Call Ventana'' se ejerce a un precio $E$ (el cual asumiremos igual a $E_1$ para alinearnos con el límite inferior de la ventana, tal como se procede analíticamente en el marco teórico) solamente si el precio del activo subyacente $S$ en el tiempo de madurez $T$ se encuentra dentro del rango $[E_1, E_2]$. Por ende, las condiciones de frontera que aplican son:

\begin{enumerate}
    \item \textbf{Condición final (Pago en madurez):}
    \begin{equation}
    C(S, T) = 
    \begin{cases} 
      S - E_1 & \text{si } E_1 \leq S \leq E_2 \\
      0 & \text{en otro caso}
    \end{cases}
    \end{equation}
    Nótese que la opción expira sin valor si $S > E_2$, lo que marca la diferencia fundamental con una opción Call Europea tradicional.
    \item \textbf{Condición en el límite inferior:}
    \begin{equation}
    C(0, t) = 0 \quad \text{cuando } S = 0
    \end{equation}
    \item \textbf{Condición asintótica en el límite superior:} A diferencia de un problema de barrera clásico que impone una restricción durante la vida útil para la variable temporal arbitraria, la opción ventana requiere un comportamiento asintótico tal que:
    \begin{equation}
    C(S, t) \rightarrow 0 \quad \text{cuando } S \rightarrow \infty
    \end{equation}
    ya que para valores extremadamente altos de $S$, es seguro que la opción terminará fuera de la ``ventana'' acordada al vencimiento.
\end{enumerate}

\subsection{Transformación a Ecuación de Difusión y Solución Analítica}
Para simplificar la EDP a la forma de la ecuación del calor de la física matemática, se emplean los cambios de variable reversibles:
\begin{equation}
C = E_1 v(x, \tau), \quad S = E_1 e^x, \quad t = T - \frac{\tau}{(\sigma^2/2)}
\end{equation}
Posteriormente, se supone que $v(x, \tau) = e^{\alpha x + \beta \tau} u(x, \tau)$, con los parámetros $\alpha = \frac{1 - k}{2}$ y $\beta = -\frac{(1 + k)^2}{4}$, donde el parámetro vinculante es $k = \frac{r}{\sigma^2/2}$.

Tras la transformación, la función $u(x, \tau)$ es solución de un problema de valor inicial basado en la ecuación de difusión:
\begin{equation}
\frac{\partial u}{\partial \tau} = \frac{\partial^2 u}{\partial x^2}
\end{equation}
cuya condición inicial asociada al pago de la opción se traduce a:
\begin{equation}
u(x, 0) = g(x) = 
\begin{cases} 
0 & \text{si } x < 0 \\ 
e^{-(\alpha-1)x} - e^{-\alpha x} & \text{si } 0 < x < \ln\left(\frac{E_2}{E_1}\right) \\ 
0 & \text{si } x > \ln\left(\frac{E_2}{E_1}\right) 
\end{cases}
\end{equation}

Desarrollando la solución de la ecuación de difusión mediante la esperanza $E_{N(x, \sqrt{2\tau})}(g)$ y el proceso de estandarización, se llega a la solución en términos de la función de distribución acumulada de probabilidad normal estándar ($F_Z$):

\begin{equation}
\begin{split}
C(S, t) &= S \left\{ F_Z \left( \frac{\ln\left(\frac{E_2}{S}\right) - \left(r + \frac{\sigma^2}{2}\right)(T - t)}{\sigma\sqrt{T - t}} \right) - F_Z \left( \frac{\ln\left(\frac{E_1}{S}\right) - \left(r + \frac{\sigma^2}{2}\right)(T - t)}{\sigma\sqrt{T - t}} \right) \right\} \\
&\quad - E_1 e^{-r(T-t)} \left\{ F_Z \left( \frac{\ln\left(\frac{E_2}{S}\right) + \left(\frac{\sigma^2}{2} - r\right)(T - t)}{\sigma\sqrt{T - t}} \right) - F_Z \left( \frac{\ln\left(\frac{E_1}{S}\right) + \left(\frac{\sigma^2}{2} - r\right)(T - t)}{\sigma\sqrt{T - t}} \right) \right\}
\end{split}
\end{equation}

\subsection{Algoritmo para Generación de la Superficie de Valoración}

Para ilustrar gráficamente la superficie $C = C(S,t)$ se emplea el siguiente marco computacional:
\begin{enumerate}
    \item Establecer los valores límite para la ventana de precios (ej. $E_1 = 90$, $E_2 = 110$) y madurez $T$.
    \item Construir una malla (\textit{meshgrid}) bidimensional computacional, definiendo un vector para el precio del subyacente $S \in [0.1, 2E_2]$ y otro para el tiempo de vida de la opción $t \in [0, T]$.
    \item Para cada par de coordenadas de la malla $(S, t)$, implementar de forma vectorizada los cálculos de los argumentos internos que alimentan a las funciones $F_Z$.
    \item Emplear rutinas especializadas de probabilidad como \texttt{scipy.stats.norm.cdf} para evaluar la función normal acumulativa y obtener el valor puntual de la opción $C(S,t)$.
    \item Utilizar herramientas gráficas tridimensionales (\textit{matplotlib 3D}) para proyectar las superficies, variando sistemáticamente los parámetros $r$ y $\sigma$ a fin de identificar sus repercusiones en el precio teórico de la opción.
\end{enumerate}

\subsection{Resultados y Análisis de la Superficie}

Para estudiar el comportamiento de la opción teórica, se implementó el algoritmo previamente descrito generando las superficies de valoración $C(S,t)$ bajo un esquema matricial de $3 \times 3$. Se consideró una opción ``Call Ventana'' con límites $E_1 = 90$ y $E_2 = 110$, y un tiempo de madurez $T = 1.0$ año. Se evaluaron tres escenarios de tasa libre de riesgo ($r$: $2\%$, $8\%$, $20\%$) cruzados con tres escenarios de volatilidad ($\sigma$: $5\%$, $20\%$, $40\%$). 

En la Figura \ref{fig:superficie} se consolidan los resultados gráficos.

\begin{figure}[htpb]
    \centering
    \includegraphics[width=\textwidth]{superficie_3x3.png}
    \caption{Superficie del valor de la opción ``Call Ventana'' $C(S,t)$ bajo diferentes combinaciones de tasa libre de riesgo ($r$) y volatilidad ($\sigma$). El dominio del precio abarca de 10 a 150 y el tiempo varía desde $t=0$ hasta cercano a la madurez $t=T$.}
    \label{fig:superficie}
\end{figure}

\subsubsection{Análisis de Resultados}

La inspección de la Figura \ref{fig:superficie} arroja perspectivas fundamentales sobre la sensibilidad de esta opción exótica:

\begin{itemize}
    \item \textbf{Condiciones de Frontera en el Tiempo (Eje $t$):} A medida que el tiempo $t$ se acerca a $T$, la superficie converge inexorablemente a la función de pago escalonada $\max(S-E_1, 0)$ pero rígidamente truncada a 0 para $S > E_2$. Esto es visible como una cresta pronunciada (la ``ventana'') cerca del borde frontal de las gráficas. En tiempos alejados de $T$ ($t \approx 0$), la superficie es mucho más suave, reflejando el valor temporal intrínseco asociado a la incertidumbre futura.
    \item \textbf{Efecto de la Tasa de Interés ($r$):} Analizando el movimiento vertical (de $r=2\%$ a $r=20\%$), se observa que el incremento en la tasa propulsa la superficie hacia valores ligeramente mayores dentro del rango de activación. Sin embargo, su influencia se ve matizada por la existencia del límite superior $E_2$: dado que el mercado penaliza (anulando el valor) si el precio subyacente crece demasiado, una tasa excesivamente alta empuja el activo fuera de la zona de ganancia más rápido.
    \item \textbf{Efecto de la Volatilidad ($\sigma$):} Analizando el movimiento horizontal (de $\sigma=5\%$ a $\sigma=40\%$), se aprecia el efecto más crítico para el diseñador de este instrumento. En bajas volatilidades, la ventana está perfectamente demarcada e inflada; la opción preserva su valor si $S$ comienza dentro del rango. A medida que la volatilidad estalla a $40\%$, la superficie se \textit{colapsa} drásticamente. Esto ocurre porque la varianza salvaje hace altamente probable que, sin importar de dónde parta $S$, el precio final termine abandonando la ventana (sea cayendo por debajo de $E_1$ o rompiendo la barrera superior $E_2$), lo que evapora el valor esperado del contrato.
\end{itemize}

\textbf{Conclusión:} A diferencia de una opción Call Europea clásica---donde una mayor volatilidad siempre beneficia al poseedor debido a las ganancias ilimitadas y las pérdidas acotadas---, en una opción ``Call Ventana'', \textbf{la alta volatilidad es destructiva}. La inclusión de la cota restrictiva superior $E_2$ invierte el perfil de riesgo: transforma la volatilidad extrema en una amenaza de expiración sin valor. Por consiguiente, esta opción exótica es ideal para inversionistas que tengan proyecciones macroeconómicas de un crecimiento de mercado moderado y fuertemente estable, y que estén dispuestos a sacrificar el potencial de ganancia ilimitada a cambio de primas (costos de adquisición iniciales) notablemente menores que las de una Call estándar.

\section{Composición y Simulación del Portafolio}

\subsection{Estrategia de Cobertura: \textit{Covered Call Ventana}}

Para demostrar la utilidad de la opción ventana en la gestión de riesgos, se propone la construcción de un portafolio basado en una estrategia modificada de \textit{Covered Call} (Llamada Cubierta). En una estrategia tradicional, un inversionista que posee el activo subyacente vende una opción Call estándar para recibir una prima; a cambio, renuncia a cualquier ganancia por encima del precio de ejercicio. Sin embargo, al utilizar una opción ``Call Ventana'', la dinámica de aseguramiento mejora notablemente en los escenarios de crecimiento extremo.

Considere un portafolio $\Pi$ compuesto en el instante $t=0$ por:
\begin{enumerate}
    \item La compra de una unidad del activo subyacente a precio $S_0$.
    \item La venta (posición corta) de una opción ``Call Ventana'' sobre el mismo activo, por la cual se recibe una prima inicial $C_0 = C(S_0, 0)$.
    \item La inversión de esta prima en un activo libre de riesgo (cuenta bancaria) que capitaliza continuamente a la tasa de interés $r$.
\end{enumerate}

El valor del portafolio en cualquier instante $t$ está dado por:
\begin{equation}
\Pi_t = S_t - C(S_t, t) + B_t
\end{equation}
donde $B_t = C_0 e^{rt}$ es el balance de la cuenta libre de riesgo. 

En el momento de la madurez $t=T$, el valor final del portafolio $\Pi_T$ adopta tres posibles comportamientos dependiendo de dónde termine el precio $S_T$:
\begin{equation}
\Pi_T = 
\begin{cases} 
S_T + C_0 e^{rT} & \text{si } S_T < E_1 \quad \text{(Amortiguación de pérdida)} \\
E_1 + C_0 e^{rT} & \text{si } E_1 \leq S_T \leq E_2 \quad \text{(Valor asegurado y estable)} \\
S_T + C_0 e^{rT} & \text{si } S_T > E_2 \quad \text{(Captura de ganancia extrema)}
\end{cases}
\end{equation}

\textbf{Ventaja teórica de seguridad:} Esta configuración demuestra que el portafolio es más seguro y eficiente que simplemente sostener el activo desnudo (\textit{unhedged}). En escenarios de caída de precios ($S_T < E_1$), la prima $C_0$ capitalizada actúa como un colchón financiero. En mercados sin tendencia fuerte ($E_1 \leq S_T \leq E_2$), el portafolio queda perfectamente blindado, asegurando un valor fijo independiente de las fluctuaciones. A diferencia de un \textit{Covered Call} tradicional, si ocurre un evento de alza extrema en el mercado ($S_T > E_2$), el inversionista no pierde el derecho a capitalizar esas enormes ganancias, ya que la opción ventana expira sin valor, liberando el activo subyacente.

\subsection{Algoritmo Computacional de Composición y Simulación}

Para verificar numéricamente la tesis de seguridad de este portafolio, se diseñó el siguiente algoritmo de simulación integral:

\begin{enumerate}
    \item \textbf{Inicialización:} Definir los parámetros del activo ($S_0, \mu, \sigma$) y de la opción ($E_1, E_2, T, r$).
    \item \textbf{Cálculo de la Prima Inicial:} Utilizar el módulo de valoración de la sección anterior para calcular rigurosamente $C_0 = C(S_0, 0)$. El capital seguro acumulado al final será $B_T = C_0 e^{rT}$.
    \item \textbf{Generación de Escenarios:} Utilizar el módulo de simulación del Movimiento Browniano Geométrico del Punto 1 para generar $M$ trayectorias (ej. $M=1000$) del activo subyacente, obteniendo un vector de posibles precios finales $S_T^{(j)}$ para $j = 1, \dots, M$.
    \item \textbf{Evaluación del Portafolio en Madurez:}
    Para cada simulación $j$, calcular el pago que se debe realizar por la opción vendida:
    \begin{equation}
    Pago^{(j)} = 
    \begin{cases} 
    S_T^{(j)} - E_1 & \text{si } E_1 \leq S_T^{(j)} \leq E_2 \\
    0 & \text{en otro caso}
    \end{cases}
    \end{equation}
    Calcular el valor final del portafolio asegurado: 
    \begin{equation}
    \Pi_T^{(j)} = S_T^{(j)} - Pago^{(j)} + B_T
    \end{equation}
    \item \textbf{Análisis Estadístico:} Comparar la distribución de $\Pi_T$ frente al activo sin cobertura ($S_T$). Calcular métricas de seguridad como la varianza (reducción del riesgo), el valor mínimo (peor escenario) y la media esperada.
    \item \textbf{Iteración de Parámetros:} Ejecutar el bucle de simulación variando sistemáticamente las tasas de interés ($r$) y las volatilidades ($\sigma$) para comprobar la persistencia de la robustez del portafolio asegurado frente a condiciones macroeconómicas cambiantes.
\end{enumerate}

\subsection{Resultados y Análisis de la Simulación}

Con el fin de validar el comportamiento de esta estrategia en diferentes escenarios económicos, se ejecutó el algoritmo computacional contrastando un mercado de baja volatilidad ($\sigma=10\%$) contra uno de alta volatilidad ($\sigma=30\%$). Se establecieron los siguientes parámetros comunes: $S_0 = 100$, $\mu=8\%$, $r=5\%$, límite inferior de la ventana $E_1 = 90$ y superior $E_2 = 110$. Se empleó $M=1000$ simulaciones para cada escenario, con un horizonte de $T=1$ año ($N=252$ pasos discretos).

Los resultados gráficos se resumen en la Figura \ref{fig:portafolio}, la cual detalla tanto la evolución temporal del portafolio (para 5 trayectorias de muestra) como el valor final alcanzado frente a la distribución estocástica subyacente.

\begin{figure}[htpb]
    \centering
    \includegraphics[width=\textwidth]{portafolio_seguro.png}
    \caption{Simulación de la estrategia \textit{Covered Call Ventana}. La fila superior muestra la evolución temporal conjunta del activo desnudo y el portafolio asegurado. La fila inferior exhibe el perfil de pago final $\Pi_T$ en madurez, superpuesto sobre el histograma de las 1000 trayectorias del activo $S_T$.}
    \label{fig:portafolio}
\end{figure}

\subsubsection{Análisis de Resultados}
La Figura \ref{fig:portafolio} corrobora contundentemente las ventajas estructurales planteadas en el diseño matemático:
\begin{itemize}
    \item \textbf{Amortiguación de caídas:} En las simulaciones donde el mercado pierde valor ($S_T < 90$), la línea negra gruesa (perfil del portafolio) se sitúa sistemáticamente por encima de la línea punteada gris (activo desnudo). La brecha vertical entre ambas líneas representa exactamente el valor de la prima cobrada $C_0 e^{rT}$. Esta protección se hace evidente en las gráficas temporales superiores, donde el portafolio inicia su camino desde un valor superior.
    \item \textbf{Estabilización en la Ventana:} Para la región de precios finales $90 \leq S_T \leq 110$, se evidencia claramente la consolidación de un ``piso'' o meseta horizontal en el perfil de pago negro. Independientemente de si el activo subyacente cierra en 95 o 105, el valor del portafolio $\Pi_T$ se mantiene totalmente estático y asegurado, bloqueando la rentabilidad.
    \item \textbf{Captura de Alta Cola (\textit{Tail Risk} Invertido):} La diferencia más notoria entre el \textit{Covered Call} estándar y nuestra estrategia con opción ventana ocurre en el extremo derecho del espectro ($S_T > 110$). Una vez cruzada la barrera superior de exclusión, la opción expira prematuramente, rompiendo la atadura del techo plano y devolviendo al inversor toda la pendiente alcista del subyacente (más el plus de la prima).
    \item \textbf{Sensibilidad a la Volatilidad:} Al observar la columna de baja volatilidad (izquierda), la densidad de la distribución (histograma rojo) está concentrada alrededor del umbral de estabilización. En este escenario predecible, vender la Call Ventana resulta altamente lucrativo. Sin embargo, en el escenario de alta volatilidad (derecha), la inmensa dispersión diluye la densidad probabilística fuera de la ventana. En consecuencia, la prima inicial recolectada $C_0$ tiende a ser minúscula, provocando que la línea del portafolio casi se traslape con la del activo desnudo y pierda su notoria capacidad de protección.
\end{itemize}

\textbf{Conclusión Final:} El portafolio instrumentado con la estrategia \textit{Covered Call Ventana} demuestra ser un mecanismo financieramente sólido para asegurar inversiones. Modifica de forma asimétrica el perfil de riesgo-rendimiento del inversor: inyecta liquidez garantizada que sirve de colchón anti-caídas, anula completamente el riesgo de variación en escenarios de crecimiento económico plano a expensas del costo de oportunidad, pero restituye soberanamente el potencial de rentabilidad en escenarios de mercado ultra-alcista. En síntesis, esta arquitectura híbrida resulta superior al \textit{Covered Call} convencional bajo la premisa de mercados lateralizados pero con sesgo esporádico fuertemente positivo (crecimientos imprevistos atípicos).

\section{Conclusiones Generales}
A lo largo de las tres fases de este proyecto, se ha articulado y comprobado la solidez de los modelos estocásticos en la ingeniería financiera:
\begin{enumerate}
    \item \textbf{Dinámica del Precio:} Mediante la discretización de Euler-Maruyama, se corroboró visualmente que la trayectoria temporal de un activo financiero está gobernada por la constante tensión entre la tendencia determinista ($\mu$) y la incertidumbre estocástica ($\sigma$).
    \item \textbf{Estructuración Exótica:} La valoración analítica y gráfica de la opción ``Call Ventana'' reveló que las cotas restrictivas ($E_1$ y $E_2$) alteran de raíz el paradigma de las opciones tradicionales. Específicamente, demostramos que este derivado exótico requiere escenarios de \textit{baja volatilidad} para maximizar su valor, ya que el exceso de dispersión amenaza con expulsar el precio final fuera de la zona de cobro, evaporando la inversión.
    \item \textbf{Gestión Inteligente de Riesgo:} La simulación del portafolio \textit{Covered Call Ventana} ratificó que es matemáticamente posible moldear el perfil de riesgo de una cartera. Al vender esta opción exótica, el inversor asegura financieramente su portafolio frente a mercados moderados o bajistas (gracias al ingreso por prima) sin hipotecar el derecho a capitalizar eventos extraordinarios al alza (gracias a la extinción prematura de la obligación).
\end{enumerate}
En definitiva, la combinación sinérgica entre modelos teóricos rigurosos (las EDP de difusión) y las herramientas computacionales avanzadas (simulaciones de Monte Carlo en Python) confiere al analista una capacidad prescriptiva invaluable para el diseño óptimo de estrategias en los mercados de capitales modernos.

\end{document}
